\subsection{The actual coefficient-tail resolution has the nonflat zero-form defect}
\label{arch136:actual-tail-syzygy}

\paragraph{Archive and input scope.}
This addresses the strict coherent Moore realization for cote136,
physical p.~117/footer116/scan117 and physical p.~121/footer120/scan121.
The bottom-module formula is the one in
ARCH-136-COHERENT-PD-REALIZATION. The closed-element calculation and
the pure divided-power detector are those of
ARCH-136-NONFLAT-TAIL-DISK, repeated below in their needed form.
The resolution, its exactness, and its minimal-retract property
are proved here; no general theorem about complete intersections
or Tate resolutions is imported. Nakayama's lemma is a
published/background input: Stacks Project, Lemma~10.20.1(2),
Tag~00DV, \url{https://stacks.math.columbia.edu/tag/00DV}.
The other arguments are tagged none, apart from the explicitly
named earlier statements.

Fix a prime \(p\), put \(R=\mathbb Z_{(p)}\) and
\(S=\Gamma_R(T)\), and let \(N=p^a\). Write
\(I_N=\bigoplus_{w\geq N}R T^{[w]}\).
For a graded \(S\)-module \(M\), its free even special PD disk is
\[
 \mathscr D_S(M)=\Gamma_S(M[0])\otimes_S\Lambda_S(M[1]),
 \qquad d x_m=e_m.                                      \tag{AT.1}
\]
Odd generators are strictly exterior, also when \(p=2\).

\paragraph{ARCH-136-ACTUAL-TAIL-SYZYGY.}
Let \(P_\bullet\to I_N\) be any weight-graded free
\(S\)-resolution whose generator weights are at least \(N\),
as permitted in the coherent construction, and let \(D\geq4\) be even. Set
\[
 P=P_{D-2},\quad Q=\ker(P_{D-2}\to P_{D-3}),\quad
 M_0=P+p^{-1}Q\subset P[1/p].                            \tag{AT.2}
\]
Then \(H^0(\mathscr D_S(M_0))\) contains nonzero elements
of exact order \(p\) in an unbounded sequence of weights.
The assertion remains true after a common positive weight shift
of the parameter modules, and after tensoring them with a
nonzero free abelian group in a fixed simplex.

In particular, the issue is present for the actual syzygy of an
allowed coefficient-tail resolution, rather than only for an
abstract module with the same saturation properties. In the
unrepaired strict realization of
ARCH-136-COHERENT-PD-REALIZATION, using a tail \(I_N\) and a
nonempty relative parameter module in some simplex, these
classes give unbounded \(p\)-torsion in vertical degree zero.
For an integral free resolution the same assertion follows
after localization at \(p\), and consequently gives nonzero
order-\(p\) classes integrally as well.

This does not contradict ARCH-136-COHERENT-LOW-REPAIR: that
statement explicitly applies the relative repair through
degrees zero and one afterwards. No conclusion about the
result of that repair, its degree-two cohomology, or the main
connectedness assertion follows here.

\paragraph{A free coefficient decomposition.}
Let \(t=T^{[N]}\) and let \(C\subset S\) be the PD
subalgebra generated by \(t\). It is isomorphic to
\(\Gamma_R(t)\). Indeed
\[
 \gamma_r(T^{[N]})=
       \frac{(rN)!}{r!(N!)^r}\,T^{[rN]},                 \tag{AT.3}
\]
and the displayed coefficient is a \(p\)-unit. Legendre's
formula gives this directly, since the base-\(p\) digits of
\(rN\) are those of \(r\) shifted by \(a\) places.
For \(0\leq s<N\), the coefficient
\(\binom{rN+s}{s}\) is likewise a \(p\)-unit. Thus
\[
 S=\bigoplus_{0\leq s<N} C\,T^{[s]},\qquad
                         I_N=S C_+.                    \tag{AT.4}
\]
The direct sum is a free \(C\)-module decomposition, with
the stated weight shifts.

Put \(x_i=\gamma_{p^i}(t)\), of weight \(h_i=Np^i\).
There are units \(\kappa_i\in R^\times\) with
\[
                 x_i^p=p\kappa_i x_{i+1}.               \tag{AT.5}
\]
The digit monomials \(\prod_i x_i^{a_i}\),
\(0\leq a_i<p\), are a unit rescaling of the divided-power
basis of \(C\). Consequently
\[
 C/pC\simeq\bigotimes_{i\geq0}
                   \mathbb F_p[x_i]/(x_i^p).             \tag{AT.6}
\]
These facts require only the same factorial valuation as (AT.3).

\paragraph{An exact minimal resolution with an elementary check.}
Use homological degree for the next complex, and denote its
differential by \(\partial\). Form
\[
 F=C\otimes_R\Lambda_R(e_0,e_1,\ldots)
                   \otimes_R\Gamma_R(f_0,f_1,\ldots),
 \quad |e_i|=1,\quad |f_i|=2,
                                                               \tag{AT.7}
\]
with weights \(\operatorname{wt}(e_i)=h_i\) and
\(\operatorname{wt}(f_i)=p h_i\), and set
\[
 \partial e_i=x_i,
 \qquad
 \partial f_i=x_i^{p-1}e_i-p\kappa_i e_{i+1}.            \tag{AT.8}
\]
Extend by the special PD derivation rule. Equation (AT.5)
gives \(\partial^2=0\) on the generators. It gives the
same identity on their divided powers because the differential
of an even generator is an odd element with square zero.
Augment \(F_0=C\) to \(R=C/C_+\).

Modulo \(p\), (AT.6) identifies this complex with the tensor
product over \(\mathbb F_p\) of the complexes
\[
 K_i=\mathbb F_p[x_i]/(x_i^p)\otimes\Lambda(e_i)
                        \otimes\Gamma(f_i),
 \qquad \partial e_i=x_i,
 \quad \partial f_i=x_i^{p-1}e_i.                        \tag{AT.9}
\]
Each \(K_i\) resolves \(\mathbb F_p\). In every positive
degree its free rank-one modules over
\(\mathbb F_p[x_i]/(x_i^p)\) have alternating differentials
\(x_i\) and \(x_i^{p-1}\). Their images are exactly the
next kernels. Equivalently, the basis pairs
\(x_i^r e_i\gamma_n(f_i)\mapsto
  x_i^{r+1}\gamma_n(f_i)\), for \(0\leq r<p-1\),
and \(\gamma_{n+1}(f_i)\mapsto
  x_i^{p-1}e_i\gamma_n(f_i)\) leave only the constant.
Tensor products over the field preserve this exactness. The
infinite tensor product is the filtered union of the finite
ones, or can be checked weight by weight, so it also resolves
\(\mathbb F_p\).

In each fixed weight, every term of \(F\) is finite free
over \(R\). Hence every homology group in that weight is
finitely generated over the noetherian ring \(R\).
The coefficient exact sequence gives an injection
\[
 H_j(F)^w/pH_j(F)^w\ \hookrightarrow\ H_j(F/pF)^w.
\]
For \(j>0\) the right side is zero; Nakayama's lemma
(Stacks Project, Lemma~10.20.1(2), Tag~00DV), with
\(pR=\operatorname{rad}(R)\), therefore gives \(H_j(F)^w=0\). Degree zero is visibly
\(R\). This proves that (AT.7)--(AT.8) is a free resolution
of \(R\) over \(C\), with no unproved acyclic-closure input.
All entries of its positive differentials lie in
\(\mathfrak n_C=(p,C_+)\).

By (AT.4), the shifted augmented complex
\[
 \mathcal P_j=S\otimes_C F_{j+1}\quad(j\geq0),
                         \mathcal P_\bullet\to I_N       \tag{AT.10}
\]
is a free \(S\)-resolution. All its generators have weight
at least \(N\), and it has finitely many generators in each
weight. All differential entries lie in
\(\mathfrak n=(p,S_+)\).

\paragraph{The detected pair in the actual odd differential image.}
First use the resolution \(\mathcal P\) in (AT.10).
Put \(r=(D-2)/2\geq1\). Fix any index \(i\), and write
\(x=x_i\), \(y=x_{i+1}\), \(h=h_i\),
\(c=p\kappa_i\), \(f=f_i\), \(e=e_i\), and
\(e'=e_{i+1}\). In \(S\otimes_C F_{2r}\), take the
two basis vectors
\[
               A=\gamma_r(f),\qquad
               B=\gamma_{r-1}(f)e e'.
\]
The two basis vectors
\(\alpha=\gamma_r(f)e\) and
\(\beta=\gamma_r(f)e'\) in
\(\mathcal P_{D-2}=S\otimes_C F_{2r+1}\) have images
\[
 u=\partial\alpha=xA+cB,
 \qquad v=\partial\beta=yA+x^{p-1}B.                    \tag{AT.11}
\]
The signs follow from
\((x^{p-1}e-c e')e=c e e'\).
Let \(L=\operatorname{im}(\mathcal P_{D-2}\to
\mathcal P_{D-3})\). In this actual syzygy,
\[
                      yu=xv,\qquad x^{p-1}u=cv.          \tag{AT.12}
\]
Thus in \(\Gamma_S(L)\) the element
\[
              K_i=y^{p-1}\gamma_p(u)-c\gamma_p(v)        \tag{AT.13}
\]
has differential zero in the free even disk, and is killed
by \(p\). For completeness, applying \(\gamma_{p-1}\)
to the first relation in (AT.12) gives
\(y^{p-1}\gamma_{p-1}(u)=
  x^{p-1}\gamma_{p-1}(v)\).
Multiplication by \(du\), followed by the differential of
the second relation, proves \(dK_i=0\). Multiplication by
\(u\) instead proves \(pK_i=0\).

This element is nonzero. The presentation
\(\mathcal P_{D-2}\to L\) has kernel
\(\operatorname{im}\mathcal P_{D-1}\). Every relation
therefore has all its coefficients in \(\mathfrak n\), by
the minimality just proved. Map the underlying coefficient ring
\(S\) to
\[
              A_p=\mathbb F_p[X,Y]/(X^p,Y^p)
\]
by retaining the digit generators represented by \(x,y\)
and killing the others. Units can be absorbed in their labels.
The ideal \(\mathfrak n\) maps into \((X,Y)\), whose
elements have zero \(p\)-th power.

In arity \(p\), projection of the free module PD algebra
onto its pure generators \(\gamma_p(\alpha)\) and
\(\gamma_p(\beta)\) descends to
\[
            \Gamma_{A_p}^p(L\otimes_S A_p)
                                      \longrightarrow A_p^2.
                                                               \tag{AT.14}
\]
Indeed products of positive arities have zero pure projection,
since their pure coefficients contain \(\binom pj=0\)
for \(0<j<p\). The pure projection of the \(p\)-th
divided power of a linear relation is the vector of \(p\)-th
powers of its coefficients, also zero here. All arity-\(p\)
relations in the generated PD ideal are covered by these two
observations. Tensor right exactness suffices to retain the
linear presentation; no flatness is used. As before, this
detector uses the underlying module PD functor and a ring map,
not a morphism preserving the coefficient PD operations.

The image of (AT.13) under (AT.14) is
\((Y^{p-1},0)\neq0\). Hence \(K_i\) has exact order
\(p\). Its weight is \(p^2(r+1)h_i\).

\paragraph{Passage to the bottom lattice.}
For \(P=\mathcal P_{D-2}\) and \(Q=\ker(P\to L)\),
let \(M_0=P+p^{-1}Q\). Since
\(x^{p-1}\alpha-c\beta\in Q\), the elements
\[
 \widetilde u=\alpha,\qquad
 \widetilde v=\beta+
          \frac{x^{p-1}\alpha-c\beta}{c}
                   =\frac{x^{p-1}}c\alpha               \tag{AT.15}
\]
belong to \(M_0\). They satisfy both relations in (AT.12).
The analogous expression to (AT.13) is closed and killed
by \(p\), and maps to the nonzero \(K_i\) under
\(M_0\to L\). It is therefore nonzero of exact order
\(p\) in \(H^0(\mathscr D_S(M_0))\).
Letting \(i\) tend to infinity makes its weights unbounded.
A common weight shift \(v\) of all parameter modules changes
these weights to \(p^2(r+1)h_i+pv\), without altering any
identity or detector. A chosen basis vector in a nonzero free
tensor factor gives a direct summand, so that assertion also
persists by inclusion and projection.

\paragraph{Independence of the chosen free resolution.}
We justify the word ``any'' in the statement. Let
\(P_\bullet\to I_N\) be another graded free resolution.
The elementary comparison construction, lifting free
generators successively through cycles and boundaries, gives
graded chain maps
\[
 \mathcal P\xrightarrow{\ i\ }P
       \xrightarrow{\ r\ }\mathcal P
\]
over the identity of \(I_N\), with
\(ri\) chain homotopic to the identity. The same successive
lifting gives a degree-preserving chain homotopy. Since all
differentials of \(\mathcal P\) have entries in
\(\mathfrak n\), every matrix of \(ri-1\) has entries
in \(\mathfrak n\).

In each homological degree the endomorphism \(ri\) is a
graded automorphism. Order the free generators by their weights.
The constant-coefficient block at any one weight is a finite
square matrix over \(R\) congruent to the identity modulo
\(p\), hence invertible. Every other coefficient maps a
generator to generators of strictly smaller weight. Induction
on weight therefore constructs an inverse: only finitely many
generator weights and finitely many generators occur below
any fixed bound. The resulting inverse is \(S\)-linear,
and it commutes with the differential because \(ri\) does.
Replacing \(r\) by \((ri)^{-1}r\) makes
\(ri=1\) strictly as chain maps.

Thus \(\mathcal P\) is a chain retract of \(P\).
The maps in degree \(D-2\) carry kernels to kernels and
extend to the lattices \(P_{D-2}+p^{-1}Q\); the composite
on the lattice belonging to \(\mathcal P\) is still the
identity. The relative module PD algebra and the even-disk
construction preserve this retract. Hence the nonzero closed
classes already exhibited survive for the bottom lattice
of every such \(P\). No finiteness assumption on the second
resolution is needed; the inverse was constructed on the
explicit finite-weight resolution \(\mathcal P\).

\paragraph{The raw strict realization, and the exact stopping point.}
In the strict presentation of
ARCH-136-COHERENT-PD-REALIZATION, the only new form-zero
relations are precisely the linear relations defining
\(M_0\), tensored with the free relative simplex module
and shifted by the native weight. The higher generators have
positive form degree. In a fixed simplex, the map from the
relative PD algebra on this form-zero module into the
underlying graded algebra of the realization has a retraction:
send the old algebra to coefficients by any vertex map, keep
the new zero-form parameters, and send all new generators of
positive form degree to zero. The linear identifying relations
are respected. This is a retraction of graded PD algebras;
no compatibility with the form differential is asserted or
needed for detecting nonzero elements.

The expressions (AT.13)--(AT.15) are closed under any special
derivation on their zero-form parameters satisfying their
linear relations, by the explicit calculation above. They
therefore remain nonzero closed order-\(p\) zero-forms in
that raw realization. There are no incoming differentials
from negative form degrees.

For a construction over \(\mathbb Z\), localization of its
free resolution and PD presentation at \(p\) gives exactly
the setting just proved. Cohomology in degree zero commutes
with this localization. A nonzero class of order \(p\) after
localization can be represented by an integral class, and
multiplied by an integer prime to \(p\) to make its order
exactly \(p\) integrally. This operation preserves its weight
and its nonzero localized image.

The low-degree repair applies precisely to such additional
closed zero-forms and may kill them while retaining all target
maps. Consequently this theorem establishes the necessity of
that step for the indicated prime-power tails; it proves no
failure of the repaired construction. In particular it does
not address whether the repaired construction can maintain
uniform vanishing in form degree two and higher.

\paragraph{Verification status.}
The author has read the complete saved proof, including the mod-\(p\)
resolution, finite-weight exactness, actual odd-image pair, minimal
presentation detector, bottom-lattice lift, strict chain retraction,
and the raw-realization and integral-localization consequences.
The parent and reviewers123 and1348 have each read the complete
saved proof and communicated PASS on 2026-11-22. These are separate
checks of this statement, not extensions of the earlier structural
syzygy checks. The subsequent Nakayama citation adds a source locator
and does not change the mathematical proof.
